% Analysis Catalog — LaTeX version of Connectivity/outputs/analysis_catalog.md
% Figures live in ./figures/ (copies from connectivity comparison outputs).
\documentclass[11pt,a4paper]{article}
\usepackage[margin=2.2cm]{geometry}
\usepackage{graphicx}
\usepackage{float}
\usepackage{hyperref}
\usepackage{xcolor}
\hypersetup{colorlinks=true, linkcolor=blue!60!black, urlcolor=blue!60!black}

\graphicspath{{figures/}}

\title{Analysis Catalog}
\author{}
\date{Generated: 2026-04-07}

\begin{document}
\maketitle

\noindent
This document is a quick guide to the analyses and figures in the connectivity comparison workflow.
It corresponds to \texttt{connectivity\_comparison/outputs/analysis\_catalog.md}.

% --- helper: full-width figure with safe scaling ---
\newcommand{\catalogfig}[2]{%
  \begin{figure}[H]
    \centering
    \includegraphics[width=\linewidth,height=0.92\textheight,keepaspectratio]{#1}
    \caption{#2}
  \end{figure}
}

\section{Matrix comparison through time}

\textbf{What it does.} Tracks whole-matrix agreement between the single estimate and the bootstrap mean across dates.
It combines correlation and error-size summaries so you can see both similarity and mismatch.

\catalogfig{time_series_matrix_metrics.png}{Whole-matrix agreement through time (correlation and error metrics).}

\textbf{How to read it.} Higher correlation means the overall structure is more similar.
Lower RMSE, MAE, and Frobenius norm mean the two products differ less in absolute terms.

\section{Threshold sensitivity}

\textbf{What it does.} Repeats key comparisons after redefining what counts as an active edge.
This checks whether the main conclusions depend on weak links or remain stable under stricter cutoffs.

\catalogfig{threshold_sensitivity.png}{Sensitivity of comparisons to connectivity thresholds.}

\textbf{How to read it.} If the curves stay similar across thresholds, the interpretation is robust.
If they change sharply, conclusions depend strongly on how small links are treated.

\section{Edge uncertainty through time}

\textbf{What it does.} Summarizes how often the single estimate falls outside bootstrap uncertainty intervals.
It also compares node-level coverage and sample-to-sample similarity through time.

\catalogfig{uncertainty_time_series.png}{Edge uncertainty and single-estimate consistency through time.}

\textbf{How to read it.} Higher outside-CI fractions mean the single estimate is less consistent with bootstrap uncertainty.
If single-vs-mean similarity is close to within-bootstrap similarity, the single matrix behaves more like a typical bootstrap realization.

\section{Block heatmap panel}

\textbf{What it does.} Shows large-scale spatial structure in the single matrix, bootstrap mean, their difference, and bootstrap uncertainty.
The matrices are block-averaged so broad patterns can be seen without plotting all edges directly.

\catalogfig{heatmap_panel_time_01.png}{Block-averaged heatmaps: single, bootstrap mean, difference, and uncertainty.}

\textbf{How to read it.} Large coherent patches mean differences are spatially structured, not just random noise.
If difference, SD, or CV hotspots line up with important regions, those regions deserve closer ecological interpretation.

\section{Edge scatter}

\textbf{What it does.} Compares sampled edge weights from the single matrix against the bootstrap mean.
The raw and log views help reveal both strong-link agreement and behavior among weaker edges.

\catalogfig{scatter_time_01.png}{Edge weights: single versus bootstrap mean (raw and log scales).}

\textbf{How to read it.} Points close to the diagonal mean good agreement.
Systematic spread away from the diagonal means one product is consistently stronger, weaker, or more variable than the other.

\section{Edge and node distributions}

\textbf{What it does.} Compares sampled edge weights plus outgoing and incoming node strengths.
This shows whether the single estimate and bootstrap mean differ in overall spread, skew, or central tendency.

\catalogfig{distributions_time_01.png}{Distributions of edge weights and node strengths.}

\textbf{How to read it.} Similar histograms suggest broadly similar strength structure.
Clear shifts or tail differences suggest the bootstrap mean is smoothing or redistributing connectivity.

\section{Top sources and sinks}

\textbf{What it does.} Compares the strongest exporters and importers under both products.
It focuses on the reefs most likely to matter for ecological interpretation and management.

\catalogfig{rank_panel_time_01.png}{Top source and sink reefs under single and bootstrap-mean products.}

\textbf{How to read it.} If the lines track closely, the most important reefs are stable.
Large mismatches suggest bootstrapping changes which reefs appear most influential.

\section{Rank correlation through time}

\textbf{What it does.} Tracks source and sink rank stability through time using the bootstrap distribution.
It summarizes whether important reefs stay important across methods and dates.

\catalogfig{rank_correlation_time_series.png}{Rank correlation of sources and sinks through time.}

\textbf{How to read it.} Higher values mean the identity of important reefs is more stable.
Drops through time suggest that bootstrapping changes ecological priorities more strongly at some dates than others.

\section{Rank-change bars}

\textbf{What it does.} Highlights reefs with the biggest changes in source or sink rank.
This makes the main movers easy to inspect without scanning the full rank table.

\catalogfig{rank_change_bars_time_01.png}{Largest rank changes for sources and sinks.}

\textbf{How to read it.} Large bars indicate reefs whose apparent importance changes the most.
Those reefs are good candidates for follow-up ecological interpretation or mapping.

\section{Profile divergence histogram}

\textbf{What it does.} Measures how much each reef's normalized destination or origin profile changes.
This asks whether reefs connect to different places even when total strength stays similar.

\catalogfig{profile_divergence_hist_time_01.png}{Distribution of profile divergence between products.}

\textbf{How to read it.} Values near zero mean the connectivity profile is similar between products.
A long right tail means a subset of reefs changes where it connects even if broad totals are stable.

\section{Profile divergence top changes}

\textbf{What it does.} Shows the reefs with the strongest destination-profile or origin-profile change.
It focuses attention on the nodes where ecological routing changes most.

\catalogfig{profile_top_changes_time_01.png}{Reefs with largest profile divergence.}

\textbf{How to read it.} Higher bars mean stronger rewiring of connectivity patterns.
These reefs may keep similar total strength while shifting their main partners.

\section{Grouping similarity}

\textbf{What it does.} Compares simple higher-level node groupings between the single matrix, bootstrap mean, and bootstrap samples.
This gives a broad system-organization view without relying on heavy community-detection methods.

\catalogfig{grouping_similarity_time_series.png}{Similarity of coarse node groupings through time.}

\textbf{How to read it.} Higher similarity means the large-scale system structure is more stable.
Lower similarity suggests bootstrapping changes how the network separates into major functional groupings.

\section{Lorenz curves}

\textbf{What it does.} Shows how concentrated total connectivity is among edges.
It is the same logic used in inequality analysis: do a few links carry most of the total connectivity?

\catalogfig{lorenz_curves_time_01.png}{Lorenz curves for connectivity concentration.}

\textbf{How to read it.} Curves farther from the diagonal mean stronger concentration in fewer links.
If the single curve is more bowed than the bootstrap mean, the single estimate is more dominated by a few pathways.

\section{Concentration through time}

\textbf{What it does.} Tracks Gini and top-share metrics through time.
This shows whether connectivity is consistently diffuse or concentrated into a small set of strong pathways.

\catalogfig{concentration_time_series.png}{Gini and top-share concentration metrics through time.}

\textbf{How to read it.} Higher values mean connectivity is more concentrated in a small number of links.
If the single estimate stays higher than the bootstrap mean, it suggests sharper or less smoothed pathways.

\section{Tail behavior through time}

\textbf{What it does.} Tracks the strongest-edge and top-k edge summaries through time.
It focuses on the upper tail where ecologically dominant pathways are most likely to appear.

\catalogfig{tail_metric_time_series.png}{Tail metrics for dominant pathways through time.}

\textbf{How to read it.} Higher tail metrics mean stronger emphasis on dominant pathways.
If only the tail changes while broad correlations stay high, differences are concentrated in the strongest links rather than everywhere.

\section{Top-edge strength distribution}

\textbf{What it does.} Compares the ranked strongest edge weights in the single matrix and bootstrap mean.
It makes it easy to see whether one product has a sharper high-end tail.

\catalogfig{top_edge_strength_distribution_time_01.png}{Distribution of ranked top edge strengths.}

\textbf{How to read it.} A heavier single tail suggests sharper pathways in the single estimate.
Closer curves suggest the strongest pathways keep similar magnitude under bootstrap averaging.

\section{Node strength uncertainty}

\textbf{What it does.} Shows bootstrap intervals for selected source and sink strengths.
It directly compares the single node value with the bootstrap range for important reefs.

\catalogfig{node_strength_uncertainty_time_01.png}{Bootstrap intervals for selected node strengths.}

\textbf{How to read it.} If the single point falls outside the interval, that node is less robust.
Wide intervals mean the reef's apparent importance is uncertain even if its mean value is high.

\section{Node CV distribution}

\textbf{What it does.} Shows how uncertain node strengths are across reefs using bootstrap coefficients of variation.
It helps separate consistently important reefs from highly unstable ones.

\catalogfig{node_strength_cv_hist_time_01.png}{Distribution of node strength coefficients of variation.}

\textbf{How to read it.} Higher CV means the node's importance is less stable.
A broad distribution means some reefs are much more uncertain than others.

\section{Top-link overlap through time}

\textbf{What it does.} Tracks how much the strongest pathways overlap between the single matrix and bootstrap mean.
This is often more ecologically meaningful than comparing all edges equally.

\catalogfig{topk_overlap_time_series.png}{Overlap of top-\(k\) links through time.}

\textbf{How to read it.} Higher overlap means the dominant pathways are more robust.
Lower overlap means bootstrapping changes which links you would identify as most important.

\section{Top-link churn}

\textbf{What it does.} Shows which edges enter or leave the strongest-link set.
It identifies specific pathways that are gained or lost when moving from the single matrix to the bootstrap mean.

\catalogfig{edge_churn_time_01.png}{Edges entering and leaving the strongest-link set.}

\textbf{How to read it.} Large churn means the identity of dominant pathways is changing.
If the entering and leaving edges have similar weights, conclusions about the strongest pathways are less stable.

\section{Stable-edge counts}

\textbf{What it does.} Counts edges that remain important across many bootstrap samples under different stability rules.
This summarizes how large the robust bootstrap-supported backbone is.

\catalogfig{stable_edge_counts.png}{Counts of stable edges under frequency rules.}

\textbf{How to read it.} More stable edges mean the bootstrap ensemble supports a clearer robust backbone.
Sharp drops under stricter frequency rules mean the strongest pathways are sensitive to sampling variation.

\section{Stable-edge overlap}

\textbf{What it does.} Compares stable bootstrap edges with the single and bootstrap-mean strongest links.
It tests whether the pathways highlighted by each product are also stable across the bootstrap ensemble.

\catalogfig{stable_edge_overlap.png}{Overlap between stable edges and single or mean strongest links.}

\textbf{How to read it.} Higher overlap means the main pathways are robust across uncertainty rules.
Low overlap means the visually strongest links may not be the most stable ones.

\section{Entropy vs strength}

\textbf{What it does.} Shows whether strong nodes are diffuse or focused in where they connect.
It links total importance to how broadly or narrowly that importance is distributed across partners.

\catalogfig{entropy_strength_time_01.png}{Node strength versus connectivity entropy.}

\textbf{How to read it.} High entropy means connections are spread across many partners.
Low entropy at high strength suggests a reef is important because of a small number of dominant routes.

\section{Dominance ratio comparison}

\textbf{What it does.} Compares how much each node is dominated by one main partner.
This helps distinguish broad exporters/importers from reefs driven by one dominant connection.

\catalogfig{dominance_ratio_time_01.png}{Dominance ratio: single versus bootstrap mean.}

\textbf{How to read it.} Higher dominance means one connection dominates the node's pattern.
If dominance falls in the bootstrap mean, bootstrapping is making connectivity appear more diffuse.

\section{Effective-link distribution}

\textbf{What it does.} Shows the effective number of destinations or sources for nodes.
It translates entropy into an easier ecological interpretation: how many meaningful partners does a reef effectively have?

\catalogfig{effective_links_distribution_time_01.png}{Distribution of effective numbers of links.}

\textbf{How to read it.} Higher values mean more diffuse connectivity patterns.
Lower values mean connectivity is concentrated into a smaller set of meaningful partners.

\section{Backbone overlap}

\textbf{What it does.} Compares a simple backbone proxy between the single and bootstrap-mean networks.
The backbone is defined transparently from thresholds and local row-share dominance rather than a complex black-box method.

\catalogfig{backbone_overlap_time_series.png}{Overlap of backbone proxies through time.}

\textbf{How to read it.} Higher overlap means both products emphasize similar core links.
Lower overlap means the structural backbone you would report depends on the uncertainty treatment.

\section{Coverage rates}

\textbf{What it does.} Summarizes whether single and bootstrap-mean system metrics fall inside bootstrap intervals.
It extends coverage beyond edges to broader quantities such as connectance, Gini, self-recruitment, and top-edge summaries.

\catalogfig{coverage_rates_time_series.png}{Coverage of system metrics relative to bootstrap intervals.}

\textbf{How to read it.} Higher outside-CI rates mean the bootstrap ensemble is adding important uncertainty information.
If the single estimate often falls outside the bootstrap range, it should not be treated as a typical realization.

\section{Coverage examples}

\textbf{What it does.} Shows example bootstrap intervals for selected node-level metrics.
It gives a concrete view of which reefs look stable and which ones move outside the bootstrap-supported range.

\catalogfig{coverage_examples_time_01.png}{Example bootstrap intervals for node-level metrics.}

\textbf{How to read it.} Points outside the interval indicate disagreement with bootstrap uncertainty.
Wide intervals mean high uncertainty even when the single point remains inside the interval.

\section{Discrepancy vs variability}

\textbf{What it does.} Compares single-vs-mean differences with the normal variability among bootstrap samples.
This is a direct check of whether the single estimate is just another plausible realization or something more distinct.

\catalogfig{single_vs_bootstrap_variability.png}{Single versus bootstrap mean relative to ensemble variability.}

\textbf{How to read it.} If the single-vs-mean value is outside the bootstrap range, the single estimate is not just a typical bootstrap realization.
If it sits inside the bootstrap range, the single matrix is more consistent with ensemble variability.

\section{Top-k probability}

\textbf{What it does.} Shows the probability that selected reefs are in the top-\(k\) set across bootstrap samples.
This turns rank stability into a direct probability of being an important source or sink.

\catalogfig{topk_probability_time_01.png}{Probability of reefs appearing in the top-\(k\) set.}

\textbf{How to read it.} Probabilities near one indicate very stable importance.
Intermediate probabilities mean a reef may be important in some realizations but not others.

\section{Rank uncertainty intervals}

\textbf{What it does.} Shows bootstrap rank intervals for important reefs.
It helps separate reefs with stable ranking from reefs whose rank changes widely across bootstrap samples.

\catalogfig{rank_uncertainty_intervals_time_01.png}{Bootstrap rank intervals for selected reefs.}

\textbf{How to read it.} Wide intervals mean reef importance is uncertain.
If the single rank falls outside the interval, the single estimate is not well aligned with bootstrap rank uncertainty.

\section{Ordination cloud}

\textbf{What it does.} Places bootstrap samples, the single estimate, and the bootstrap mean in low-dimensional space.
This gives a compact view of whether the single matrix sits inside or outside the bootstrap cloud.

\catalogfig{ordination_time_01.png}{Low-dimensional ordination of bootstrap samples and single estimate.}

\textbf{How to read it.} If the single estimate sits inside the bootstrap cloud, it is more consistent with the ensemble.
If it sits apart from the cloud, bootstrapping is adding information not captured by the single estimate.

\section{Recommended figures for a paper or report}

\begin{itemize}
  \item \texttt{time\_series\_matrix\_metrics.png} for the broad comparison.
  \item \texttt{rank\_correlation\_time\_series.png} and\\ \texttt{rank\_change\_bars\_time\_XX.png} for source and sink stability.
  \item \texttt{lorenz\_curves\_time\_XX.png} and \texttt{concentration\_time\_series.png} for concentration of connectivity.
  \item \texttt{stable\_edge\_overlap.png} and \texttt{single\_vs\_bootstrap\_variability.png} for robustness and uncertainty.
  \item \texttt{ordination\_time\_XX.png} for the bootstrap cloud view.
\end{itemize}

\noindent
\textit{Note:} filenames with \texttt{time\_XX} refer to time indices selected by the workflow (e.g.\ \texttt{time\_01}); replace with the indices used in your run when citing figures.

\end{document}
